\documentclass[11pt]{article}
\usepackage[a4paper,margin=1in]{geometry}
\usepackage{fontspec}
\usepackage[boldfont=false]{xeCJK}
\setCJKmainfont{FandolSong-Regular.otf}
\usepackage{amsmath}
\usepackage{amssymb}
\usepackage{array}
\usepackage{booktabs}
\usepackage{graphicx}
\usepackage{adjustbox}
\usepackage{enumitem}
\setlist[itemize]{topsep=2pt,partopsep=0pt,parsep=0pt,itemsep=2pt,leftmargin=1.4em}
\setlist[enumerate]{topsep=2pt,partopsep=0pt,parsep=0pt,itemsep=2pt,leftmargin=1.6em}
\usepackage{parskip}
\usepackage{fvextra}
\fvset{breaklines=true,breakanywhere=true,fontsize=\small}
\setlength{\parindent}{0pt}

\begin{document}

\begin{center}
  {\LARGE\bfseries Analytical techniques}\\[0.35em]
  {\large A Level Chemistry}
\end{center}
\vspace{0.6em}

\section*{Thin-layer chromatography}

\textbf{Chromatography} 色谱 separates a mixture using two "phases" — one that stays still and one that moves. In \textbf{thin-layer chromatography} 薄层色谱 (TLC):

\begin{itemize}
\item the \textbf{stationary phase} 固定相 stays still (for example aluminium oxide on a plate).

\item the \textbf{mobile phase} 流动相 moves (a polar or non-polar solvent that travels up the plate).

\item the \textbf{baseline} 基线 is the starting line where the spots are placed; the \textbf{solvent front} 溶剂前沿 is the highest level the solvent reaches.

\end{itemize}

The $R_{\text{f}}$ value compares how far a spot moves with how far the solvent moves:

\[
R_{\text{f}} = \frac{\text{distance moved by the spot}}{\text{distance moved by the solvent front}}
\]

It is always between 0 and 1. A substance that sticks more strongly to the stationary phase, or is less soluble in the mobile phase, moves less and has a \textbf{smaller} $R_{\text{f}}$.

\section*{Gas/liquid chromatography}

In \textbf{gas/liquid chromatography} 气液色谱 (GLC):

\begin{itemize}
\item the stationary phase is a high-boiling-point non-polar liquid on a solid support.

\item the mobile phase is an unreactive carrier gas.

\item the \textbf{retention time} 保留时间 is how long a component takes to pass through.

\end{itemize}

The area of each peak gives the percentage of that component in the mixture. A component that interacts more with the stationary phase has a longer retention time.

\section*{Carbon-13 NMR spectroscopy}

\textbf{NMR} 核磁共振 (nuclear magnetic resonance) studies how certain nuclei behave in a strong magnetic field.

In carbon-13 NMR, each different \textbf{chemical environment} 化学环境 of carbon gives one peak. So:

\begin{itemize}
\item the \textbf{number of peaks} tells you how many different carbon environments there are (equivalent carbons share a peak).

\item the position of each peak (its chemical shift) suggests the type of carbon, which helps you deduce possible structures.

\end{itemize}

\section*{Proton (¹H) NMR spectroscopy}

Proton NMR looks at the hydrogen atoms (\textbf{proton} 质子 nuclei). From the spectrum you read off:

\begin{itemize}
\item \textbf{chemical environments}: protons in different environments appear at different \textbf{chemical shift} 化学位移 values.

\item \textbf{relative numbers}: the relative peak \textbf{areas} give the ratio of each type of proton.

\item \textbf{splitting} 裂分: a peak is split by the protons on the neighbouring carbon, following the \textbf{$n+1$ rule} — $n$ equivalent neighbours split a peak into $n+1$ lines:

\end{itemize}

\begin{center}
\adjustbox{max width=\linewidth}{%
\begin{tabular}{ll}
\toprule
\textbf{Neighbours ($n$)} & \textbf{Pattern} \\
\midrule
0 & \textbf{singlet} 单峰 \\
1 & \textbf{doublet} 双峰 \\
2 & \textbf{triplet} 三峰 \\
3 & \textbf{quartet} 四峰 \\
many & \textbf{multiplet} 多重峰 \\
\bottomrule
\end{tabular}}
\end{center}

\subsection*{Practical points}

\begin{itemize}
\item \textbf{tetramethylsilane} 四甲基硅烷 (TMS) is the standard, set at a chemical shift of $0$.

\item a \textbf{deuterated solvent} 氘代溶剂 (such as $\text{CDCl}_3$) is used so that the solvent itself gives no proton signal.

\item shaking the sample with $\text{D}_2\text{O}$ makes the O–H and N–H peaks disappear (their hydrogen is swapped for deuterium), which identifies those protons.

\end{itemize}


\end{document}
